\subsection{Results used from earlier papers}\label{sec:importados}
In this list \(m\) is the residual size and \(n\) the initial board size. On a reachable ordinary
residual board, \(p\) denotes the interior Bethe stationary point on its support, with mass one in each
row and column. We write \(\ell_f=\sum_{e\in f}p_e\), \(h_f=1-\ell_f\) for both oblique families. The
common bounds are
\begin{equation}\label{eq:clase-importada}
 \max_e p_e\le C/m,\qquad \min_f h_f\ge\eta>0,\qquad
 \lambda_2(\mathcal L_p)\ge\gamma>0 .
\end{equation}
The bipartite graph on rows and columns is connected. Changing the sign of the column coordinates
transforms its Laplacian \(\mathcal L_p\) into the axial incidence matrix weighted by \(p\), with kernel
\((\mathbf1_R,-\mathbf1_C)\). Thus the gaps in~\cite{pc} and~\cite{tc} are the same. In formulas using
\(\tau_e=p_e(1-p_e)/(1-4p_e)\), we have \(p_e\le\tau_e\le2p_e\) for large \(m\). Comparing the forms
preserves a uniform gap. The predictor uses the axial cofactor, the local factor, and the normalization
term \(\tfrac12\log(2\pi)\).

\paragraph{Child insertion and row sum~\cite[Thm.~7.1 and Cor.~7.2]{pc}.} If \(m\ge\delta_0n\)
and~\eqref{eq:clase-importada} holds, there exist \(C_*,m_0\) depending only on
\(\delta_0,C,\eta,\gamma\). For \(m\ge m_0\), every legal child has a unique stationary point on the
branch \(p<1/4\) and
\[
 L(B+e)-L(B)=\log[p_e(1+\delta_e)]+r_e,\qquad
 |\delta_e|\le C_*/m,\quad |r_e|\le C_*/m^2,\quad \sum_{e\in i}p_e\delta_e=0 .
\]
The correction \(\delta_e\) is the projected derivative of the logarithm of the prefactor divided by
\(p_e\). The corollary gives \(|\log\rho(B,i)|\le C_*/m^2\) for every row \(i\), summing over all its
children. These two estimates enter the history and insertion identities in Appendices~\ref{app:signed}
and~\ref{app:interior}.

\paragraph{Mesoscopic class~\cite[\S8.1 and Thm.~8.1]{pc}.} In the same geometry, suppose
that~\eqref{eq:clase-importada} holds and that \(\max_f\ell_f\le\epsilon\le\epsilon_0\). The thresholds
\(\epsilon_0,m_0\) and the constant \(C_*\) depend on \(C,\eta,\gamma\), independently of \(n/m\) and of
the number of lines. For \(m\ge m_0\) all legal children have their small stationary point and
\(|\log\rho(B,i)|\le C_*/m^2\). The proof retains the preceding representation for each child, giving
\(\nu(e\mid i)=p_e(1+O(m^{-1}))\) even in the transition at the first exit. This is used in
Appendices~\ref{app:signed}, \ref{app:arrival}, and~\ref{app:interior}.

\paragraph{Prefactor expansion~\cite[\S\S7.5--7.8 and 8.4]{pc}.} On a fixed support in the macroscopic
class, with the gauge above, \(s_{2,a}=\sum_{i\in a}p_i^2\), and \(\kappa_i=a_i^{\mathsf T}S_0^+a_i-4\),
we use
\[
 \ell+\tfrac12\log\det S_{0,\rm gf}+3\sum_i p_i^2
   =-14\sum_i p_i^3-\tfrac32\sum_i p_i^2\kappa_i
       +\tfrac94\sum_a s_{2,a}^2+O(m^{-2}).
\]
The constants in the remainder are class constants, depending on \(\delta_0\) in §7. The mass bounds in
§8.4 make them independent of \(n/m\) in the mesoscopic class. This is used in
Appendix~\ref{app:interior} to identify \(mV\).

\paragraph{Stopping and macroscopic persistence~\cite[Def.~2.2, \S4 and Cor.~4.12]{tc}.} Fix
\(M=\lceil\delta n\rceil\) and a floor \(K_0\le K\le M\). Up to \(M\) the discrepancy in intervals,
fibers, and rectangles must remain below the macroscopic radius. At \(M\) one also imposes the entry
tolerances of §5.3 and the maximum length \(D_{\max}(M)<\alpha_{\rm hist}M\) for empty crossings. After
that one requires the inner class of~\eqref{eq:clase-importada} with \(\ell_f\le\epsilon\), up to and
including \(K\). The first violation stops the history. For a fixed fraction \(\delta>0\), a sufficiently
small radius, a sufficiently fine partition, and large \(n\), the probability under \(\nu\) of reaching
the radius by \(M\) is \(\le C_{\delta,r}e^{-c_{\delta,r}n}\). On the complementary event the Bethe
maximizer is unique and satisfies~\eqref{eq:clase-importada} uniformly over lines. Section~4.3, (4.3),
gives \(mp_e/W(e)=1+O_\delta(\zeta+1/m)\) at available sites and the same absolute error bound for the
diagonal masses relative to \(1-u_f/U_f\), where \(\zeta\) sums the interval and fiber discrepancies, the
partition step, and the discretization error. For a fixed number of lines, Lemma~4.10 gives the one-step
hazard \(m^{-1}(\#\text{axes}+\sum_f(1-u_f/U_f))+O_\delta((\zeta+1/m)/m)\) under the law that chooses
within the row according to \(p\). The child ratio above preserves that accuracy under \(\nu\). We use
the corollary in Appendix~\ref{app:signed} and tracking at a fixed auxiliary cut-off in
Appendices~\ref{app:interior} and~\ref{app:transfer}.

\paragraph{Entry, prescriptions, and bands~\cite[\S5, Lemmas 5.2, 5.7 and Prop.~5.6, 5.8]{tc}.} Fix the
inner and outer margins with \(2\epsilon=\epsilon_{\rm out}\le\epsilon_0(C_{\rm out},\eta_{\rm
out},\gamma_{\rm out})\), and the tolerances as in §5.4. The entry value \(\delta\) must be sufficiently
small for those tolerances. Lemma~5.7 provides entry of the observables and a geometric exception \(\le
C_\lambda e^{-\lambda M}\) for \(D_{\max}(M)\). Proposition~5.6 reconstructs the inner class at \(M\) and
on boards satisfying the tolerances. Every child of a regular parent, including the exit child, inherits
the outer class. The weaker margins hold there, and the diagonal mass satisfies \(\le2\epsilon\) in the
mesoscopic regime. Thus the parents belong to the classes of the two insertion theorems, with \(m>M\) and
\(m\le M\), respectively. Under the law \(\nu\), the probability that \(s\) prescribed rows and \(t\)
prescribed columns survive until \(m\), with regular parents up to the arrival step, is
\(\le(m/n)^{s+t}\exp(C(st+s+t)/m)\), uniformly in \(s,t\). For every \(K_0\le K\le M\) and \(u=2^qK\le
M\), with \(q\ge0\) an integer, regular entry followed by first exit in \([u,2u)\) costs \(\le C
e^{-c_*u}\), with \(c_*>2\beta\), \(\beta=2\epsilon_{\rm out}/(1-\epsilon_{\rm out})\). The band
constants are common to all floors. These are the inputs about histories for Appendices~\ref{app:signed}
and~\ref{app:arrival}.

\paragraph{Static cost~\cite[Lemmas 6.5--6.6, Prop.~6.7 and Cor.~6.8]{tc}.} Let \(\widehat W_e=p_e(1-p_e)\),
and let \(\mu_{\widehat W}\) be the law on permutations supported on the mask, with weights
\(\prod_{e\in\sigma}\widehat W_e\). For an interior stationary point with unit axial masses and positive
slacks, \cite[Lemma~6.5]{tc} gives the identity
\[
 Q(B)e^{-H_B(p)}=\frac{\operatorname{per}\widehat W}{\prod_e(1-p_e)}
                  \mathbb E_{\mu_{\widehat W}}\mathcal A,\qquad
 \mathcal A(\sigma)=
 \frac{\prod_f h_f}{\prod_e(1-p_e)^2}\,\mathbf1_{\{\sigma\ {\rm legal}\}}
 \prod_{e\in\sigma}\frac{(1-p_e)^2}{h_{D(e)}h_{A(e)}} .
\]
Here \(\Gamma_4,\Gamma_2\) denote the canonical queen and rook normalizers with the weights \(p\) of the
same board. If \(Q(B)>0\), \(p_e\le C/m\), and \(\|p-\mathbf1\mathbf1^{\mathsf T}/m\|_{\rm
op}\le\rho_*<1\), one obtains \(\log Q(B)-L(B)=\log\mathbb E_{\mu_{\widehat W}}\mathcal A
-\log(\Gamma_4/\Gamma_2)+o(1)\), uniformly for fixed \(C,\rho_*\). The axial matrix has spectrum
\(1\pm\sigma_i(p)\), so \eqref{eq:clase-importada} implies the last hypothesis.

With the margins of~\eqref{eq:clase-importada}, \cite[Lemma~6.6]{tc} gives
\((S(t)^+)_{uv}=\mathbf1_{\{u=v\}}+O_{C,\eta,\gamma}(m^{-1})\), uniformly in the axial indices and in
\(0\le t\le1\). Here \(S(t)=(1-t)S_{R,2}+tS_{R,4}\), where \(S_{R,d}\) is the axial incidence matrix
weighted by \(p(1-p)/(1-dp)\). With \(\max_f\ell_f\le\epsilon_{\rm ext}<1\), \cite[Prop.~6.7]{tc} uses this
inverse to give \(\log(\Gamma_4/\Gamma_2)=O_{C,\eta,\gamma}(\epsilon_{\rm ext}+m^{-1})\). At a first
mesoscopic exit, \cite[Cor.~6.8]{tc} yields \((\log Q-L)^+\le C_0+\beta m\), where the left-hand side is
read as zero when \(Q=0\), uniformly in the outer class. Its \emph{proof} applies at every arrival in that
class. Indeed, \(\sum_f\log h_f-2\sum_e\log(1-p_e)\le O_C(1)\), since \(\sum_f\ell_f=2m=2\sum_e p_e\) and
\(\sum_ep_e^2\le C\). Each of the \(m\) chosen squares contributes at most \(2\epsilon_{\rm
ext}/(1-\epsilon_{\rm ext})\) to the logarithm of the last product. \cite[Lemmas~6.5--6.6 and Prop.~6.7]{tc}
absorb the remainder into \(C_0\), independently of \(n/m\). Appendix~\ref{app:signed} applies the corollary
in the bands, and Appendix~\ref{app:arrival} applies the static argument at regular arrival with a long
line.

